\section{How to read the formal companion}
\label{sec:reading}
The Lean project formalizes the explicitly scoped mathematical core of Chapter VI. It turns an equation such as $Ce^{tA}=e^{tB}C$ into a statement with named spaces, scalar fields, operator types, and hypotheses. This guide follows those interfaces. It explains how the formal pieces fit together and where the traditional edition goes beyond them.

The latest core is release \code{0.5.0-checked-fifth-block}. Its preserved ledger records all 33 scoped results as complete. The source contract names 341 explicit theorem declarations across 47 production modules. The preserved September 6 compiler and namespace audit reports record that 782 declarations were inspected in the project namespace. These counts measure different things: manuscript endpoints, named theorem declarations, files, and the wider namespace. None is a count of independent new mathematical discoveries.

The traditional monograph preserves the larger written development in Parts I through III, followed by the worked atlas in Part IV. Part I is the latest reviewed core. Part II retains the strengthened exposition. Part III retains earlier results not subsumed by that exposition and the bounded Banach operator material migrated from Chapter IV. The theorem-like result records across these parts include conceptual overlap. Formal completion remains attached to the 33 specific Part I ledger entries.

\begin{readerroadmap}[title={Read one endpoint through its whole interface}]
Start with the statement and proof in the traditional volume. Locate its numbered row in the 33-result ledger. Read the surrounding variables in the original Lean module, then compare the declaration's hypotheses and conclusion with the mathematical statement. Finally, identify whether the evidence is an inherited kernel record, a current source reconciliation, or a bounded computation. These steps answer different questions.
\end{readerroadmap}

\subsection{Three forms of evidence}
\begin{longtable}{@{}p{.23\linewidth}p{.70\linewidth}@{}}
\toprule Evidence & What it establishes here\\\midrule\endhead
Inherited kernel evidence & The saved September 6 build and axiom audit for the pinned source and dependencies. No new Lean compilation occurred in this edition.\\
Recorded source and computational checks & Source hashes, target inventories, verifier controls, and available finite examples were rerun in the recorded monograph review. These are not kernel checking or universal proofs.\\
Mathematical correspondence & The relationship between prose, hypotheses, definitions, ordinary real/complex specializations, and theorem interfaces is a manual review responsibility.\\
\bottomrule\end{longtable}

The source report explicitly records \code{full_paper_formalized: false}. That remains correct even though \code{all_scoped_results_complete: true} is recorded for the historical kernel run. Figures, explanatory prose, broader extensions, and three open conjectures are not swallowed by the scoped denominator.

\subsection{A practical reading route}
Read the matching traditional argument first, use the visual guide to inspect a finite example, and then find its ledger entry in Section~\ref{sec:ledger}. Follow the named endpoint to a source excerpt or the complete declaration catalog. Check the ambient variables at the top of the original module before attempting to reuse an excerpt. A snippet does not repeat every \code{variable}, \code{open}, instance, or imported notation in its file.

All displayed Lean blocks are exact contiguous source excerpts from the preserved production modules. Raw UTF-8 copies and source hashes accompany the editable guide. Font fallback in the PDF changes rendering only. The snippets are reading excerpts, not independent programs that can be compiled without their original context.

\subsection{The notation that matters most}
\begin{longtable}{@{}p{.31\linewidth}p{.62\linewidth}@{}}
\toprule Lean notation & Mathematical reading\\\midrule\endhead
\code{Matrix I J ℂ} & A complex matrix with rows indexed by $I$ and columns by $J$. Different index types permit rectangular maps.\\
\code{E →L[ℂ] F} & A bounded complex linear map from $E$ to $F$. Completeness and inner-product assumptions come from ambient instances.\\
\code{C.comp A} & Composition $C\circ A$: apply $A$ first, then $C$.\\
\code{Module.End ℂ V} & A complex linear endomorphism of $V$. Finite dimension is a separate hypothesis.\\
\code{IntegrableOn f (Ioi 0)} & Bochner integrability of $f$ over $(0,\infty)$ with the stated measure, here ordinarily Lebesgue measure.\\
\code{Tendsto f atTop} \texttt{(}\CodeSymbol{1D4DD}\texttt{ 0)} & The limit of $f(t)$ is zero as $t$ tends to $+\infty$.\\
\code{0 ≤ A} & Positivity in the operator or $C^*$-algebra order in the fractional modules. This is not entrywise matrix positivity.\\
\code{∀ i j, 0 ≤ K i j} & Entrywise nonnegativity of a finite real matrix.\\
\code{Ring.inverse} & A total inverse operation. The surrounding proof must establish the required invertibility before an inverse formula has its intended analytic meaning.\\
\bottomrule\end{longtable}

\subsection{Common mistakes that the interfaces prevent}
Column-stochastic transition matrices act on probability columns. Observable propagation uses the transpose. Probability aggregation is a rectangular mass-summing map. A weighted conditional expectation is a square projector on observables. Their identities are related but their types differ. Similarly, $e^{tA}$ describes a general generator, while a positive fractional operator produces $e^{-tA^\alpha}$. The minus sign changes the Duhamel defect in the fractional semigroup theorem.

Zero evidence, one coarse packet, a zero-dimensional space, a zero spectral value, and a nonnormal hidden generator are not interchangeable edge cases. Each belongs to a particular interface. This guide identifies the relevant hypothesis at the point where it is used.
